\section{ماژول kvtc\_compressor (KV Cache Transform Coding Compressor)}
\label{sec:module_kvtc_compressor}

\subsection{معرفی و انگیزه}
در مدل‌های زبانی بزرگ، کش کلید-مقدار (\lr{KV Cache}) به عنوان یکی از گلوگاه‌های اصلی مصرف حافظه در طول استنتاج متوالی با طول بافت بلند شناخته می‌شود. روش \lr{KV Cache Transform Coding (KVTC)} با الهام از تئوری کدگذاری تبدیل در پردازش سیگنال و تصویر، با هدف فشرده‌سازی حداکثری بازنمایی‌های فعال‌سازی حافظه کش و وزن‌های متناظر پیشنهاد شده است.

در این پروژه، علاوه بر ماژول مستقل نگهداری پویا \lr{KVTCStore} برای کدگذاری تبدیل برخط کش‌های کلید و مقدار در حین استنتاج، موتور فشرده‌ساز \lr{KVTCCompressor} به عنوان یک استراتژی کوانتایش صحیح ماتریس‌های وزن لایه‌های خطی به کار گرفته می‌شود. در این مکانیزم، داده‌ها پس از ناهمبسته‌سازی آماری به وسیله تبدیل تحلیل مؤلفه‌های اصلی (\lr{PCA})، تحت کوانتایش اسکالر وفقی قرار گرفته و با لایه‌های صحیح کم‌حجم (\lr{Int8/Int16}) به همراه شاخه‌های انطباقی \lr{LoRA} جایگزین می‌شوند.

\subsection{مبانی نظری و فرمول‌بندی ریاضی}
ایده اصلی کدگذاری تبدیل (\lr{Transform Coding})، نگاشت بردارهای همبسته به زیرفضای مؤلفه‌های اصلی ارتوگونال است که در آن مؤلفه‌ها از یکدیگر ناهمبسته شده و انرژی سیگنال در ابعاد اولیه متمرکز می‌گردد. برای تانسور داده‌های ورودی کالیبراسیون $\mathbf{X} \in \mathbb{R}^{N \times D}$ با ماتریس کوواریانس نمونه:
\begin{equation}
\mathbf{\Sigma} = \frac{1}{N} (\mathbf{X} - \mathbf{\mu})^\top (\mathbf{X} - \mathbf{\mu})
\end{equation}
تجزیه مقادیر ویژه یا تجزیه مقادیر منفرد بر روی داده‌های میانگین‌زدوده اجرا می‌شود:
\begin{equation}
\mathbf{X} - \mathbf{\mu} = \mathbf{U} \mathbf{S} \mathbf{V}^\top \implies \mathbf{P} = \mathbf{V} \in \mathbb{R}^{D \times D}
\end{equation}
که در آن $\mathbf{P}$ پایه تبدیل \lr{PCA} را تشکیل می‌دهد. داده‌های تبدیل‌یافته به صورت زیر تصویر می‌شوند:
\begin{equation}
\mathbf{Z} = (\mathbf{X} - \mathbf{\mu}) \mathbf{P}
\end{equation}
سپس کوانتایش یکنواخت اسکالر بر اساس بودجه بیتی $b$ بر روی مقادیر تصویرشده اعمال می‌گردد:
\begin{equation}
\Delta = \frac{\max(|\mathbf{Z}|)}{2^{b-1} - 1}, \quad \mathbf{Q} = \text{clip}\left( \left\lfloor \frac{\mathbf{Z}}{\Delta} \right\rceil, -2^{b-1}+1, 2^{b-1}-1 \right)
\end{equation}
آنتروپی شانون مقادیر کوانتایزشده به عنوان حد پایین تئوریک نرخ فشرده‌سازی سنجیده می‌شود:
\begin{equation}
\mathcal{H}(\mathbf{Q}) = -\sum_{s \in \mathcal{S}} p(s) \log_2 p(s)
\end{equation}
در پیاده‌سازی لایه‌های وزن خطی \lr{KVTCQuantizedLoRALinear}، ماتریس وزن لایه پایه با یک بافر عدد صحیح به همراه ضریب مقیاس (\lr{Scale}) و نقطه صفر (\lr{Zero Point}) ذخیره شده و خروجی با ترکیب خروجی عدد صحیح دکوانتایزشده و مسیر کم‌رتبه به صورت زیر محاسبه می‌شود:
\begin{equation}
\mathbf{y} = \mathbf{x} \left( \mathbf{W}_{\text{int}} \cdot s + z_0 \right)^\top + \frac{\alpha}{r} (\mathbf{x} \mathbf{A}_{\text{lora}}^\top) \mathbf{B}_{\text{lora}}^\top + \mathbf{b}
\end{equation}

\subsection{معماری و پیاده‌سازی}
ماژول در فایل \lr{Codes/framework/compression\_engine/kvtc\_compressor.py} مستقر بوده و اجزای اصلی آن عبارتند از:
\begin{itemize}
    \item \textbf{کلاس \lr{SimpleEntropyCoder}:} ابزار تخمین تجربی آنتروپی شانون بر مبنای فراوانی مقادیر گسسته برای برآورد نرخ فشرده‌سازی تئوریک.
    \item \textbf{کلاس \lr{KVTCStore}:} ذخیره‌ساز برخط تبدیل \lr{KV Cache} که ماتریس‌های پایه \lr{PCA} را برای کش‌های کلید و مقدار در لایه‌های مدل کالیبره کرده و فشرده‌سازی حافظه توالی‌های پویا را مدیریت می‌کند.
    \item \textbf{کلاس \lr{KVTCQuantizedLoRALinear}:} لایه سفارشی جایگزین لایه‌های خطی که وزن‌های اصلی را در قالب بافرهای عدد صحیح \lr{int8} یا \lr{int16} ذخیره کرده و از متغیرهای شناور \lr{scale} و \lr{zero\_point} در کنار آداپتورهای باقیمانده \lr{LoRA} برای عبور پیشرو دقیق استفاده می‌نماید.
    \item \textbf{کلاس \lr{KVTCCompressor}:} مدیریت اسکن ماژول‌های شبکه، استخراج پارامترهای لایه‌ها و ثبت در رجیستری فشرده‌سازها.
\end{itemize}

\subsection{نتایج تجربی}
ارزیابی ماژول \lr{kvtc\_compressor} بر روی مدل‌های زبانی و انکودرهای پایه بر اساس داده‌های تجربی تاییدشده در \lr{thesis\_final/master\_verified\_results.json} انجام شده و در جدول \ref{tab:benchmark_kvtc_compressor} گزارش شده است:

\begin{table}[H]
\centering
\small
\caption{نتایج تجربی و اعتبارسنجی بنچ‌مارک فشرده‌ساز \lr{kvtc\_compressor} روی مدل‌های پایه}
\label{tab:benchmark_kvtc_compressor}
\begin{tabular}{lccccc}
\toprule
\textbf{مدل پایه} \lr{(Base Model)} & \textbf{دقت} \lr{(Accuracy)} & \textbf{شاخص} \lr{(F1-Score)} & \textbf{حافظه} \lr{(Peak VRAM)} & \textbf{توان} \lr{(Throughput)} & \textbf{وضعیت} \\
\midrule
\lr{DistilBERT} & \lr{89.68\%} & \lr{0.9004} & \lr{388.11 MB} & \lr{151,447.23 tok/s} & \lr{PASS} \\
\lr{RoBERTa} & \lr{50.92\%} & \lr{0.6748} & \lr{457.87 MB} & \lr{97,731.31 tok/s} & \lr{PASS\_LOW} \\
\lr{GPT-2} & \lr{52.75\%} & \lr{0.6756} & \lr{666.42 MB} & \lr{66,198.67 tok/s} & \lr{PASS\_LOW} \\
\lr{Qwen} & \lr{49.08\%} & \lr{---} & \lr{---} & \lr{---} & \lr{PASS} \\
\bottomrule
\end{tabular}
\end{table}

بررسی داده‌های حاصل از ارزیابی استاندارد بیانگر موفقیت کامل ماژول در مدل \lr{DistilBERT} با دقت \lr{89.68\%}، شاخص \lr{F1} برابر با \lr{0.9004} و توان پردازشی بسیار چشمگیر \lr{151,447.23 tok/s} با مصرف حافظه بهینه \lr{388.11 MB} است. در مدل‌های عمیق‌تر نظیر \lr{RoBERTa} و مدل دیکودر \lr{GPT-2}، دقت در سطح \lr{50.92\%} و \lr{52.75\%} قرار گرفته و وضعیت \lr{PASS\_LOW} ثبت گردیده است. همچنین برای مدل بزرگ زبانی \lr{Qwen} با ثبت دقت \lr{49.08\%} وضعیت قبولی آزمون تایید شده است.

\subsection{محدودیت‌ها و تحلیل شکست}
تحلیل فنی ساختار ماژول و رفتار تجربی آن بر روی تسک‌های متوالی موارد زیر را نمایان می‌سازد:
\begin{itemize}
    \item \textbf{انحراف آماری در عدم تطابق داده‌های کالیبراسیون (\lr{Distributional Shift}):} پایه‌های تصویر \lr{PCA} به شدت وابسته به نمونه‌های مرحله کالیبراسیون هستند. در صورتی که توزیع متن‌های آزمون نسبت به متون کالیبراسیون تفاوت فاحشی داشته باشد، ناهمبستگی بهینه از بین رفته و نویز کوانتایش در ابعاد فرعی تقویت می‌شود که علت اصلی افت دقت در مدل‌های \lr{RoBERTa} و \lr{GPT-2} است.
    \item \textbf{سربار دکوانتایز پویا در فقدان کرنل‌های همجوش:} لایه \lr{KVTCQuantizedLoRALinear} در طول گذر پیشرو وزن‌های صحیح را در لحظه به فرمت ممیز شناور تبدیل می‌کند؛ بدون بهره‌گیری از هسته‌های سفارشی \lr{W8A16 CUTLASS}، این فرآیند موجب صرف چرخه‌های اضافی حافظه شده و پتانسیل بیشینه شتاب را محدود می‌سازد.
    \item \textbf{عدم ذخیره‌سازی واقعی آنتروپی (\lr{Entropy Bitstream Omission}):} اگرچه برآوردگر آنتروپی، حد فشرده‌سازی بسیار بالایی (تا بیش از ۸ برابر) را نشان می‌دهد، اما به دلیل عدم پیاده‌سازی کدگذار دودویی حقیقی هافمن یا کدگذاری حسابی، فشرده‌سازی فیزیکی به بایت‌های کامل محدود می‌ماند.
\end{itemize}
